% =====================================================================
%  SUBSECTION: The associator debt and the Higgs mechanism
%  Drafted 2026-07-03 from the mass-programme audit of the same date.
%
%  PLACEMENT: this is the closing subsection of the boson-masses section.
%  Either paste it at the END of boson_masses_section_2026-06-08.tex
%  (after the op:vev note that closes "The top Yukawa"), or \input this
%  file in main.tex immediately after the boson_masses_section line, in
%  which case the subsection attaches to that section's numbering.
%
%  WHAT IT DOES (decisions from the 2026-07-03 discussion):
%   - Identifies the framework's dimensionless geometric factors with the
%     Yukawa couplings: every framework mass formula m = (geometry) x v
%     IS the Yukawa relation; y_t = 1 was already a derived Yukawa.
%   - Recovers, not rivals, the mechanism: debt measured relative to the
%     selected direction; symmetric phase (no selection) => massless;
%     "why is a direction selected" == "why does EWSB occur".
%   - kappa_f = kappa_V = 1 as an automatic consequence of linearity in v,
%     grounded in the ATLAS/CMS coupling-vs-mass measurements; stated as
%     a null-forever prediction (companion snippet for sec:predictions at
%     the bottom of this file).
%   - Renames op:sector-scales in Higgs language: it IS the Yukawa-
%     normalisation problem (reframing only; the op itself is unchanged
%     in sec:masses).
%   - Makes the Lambda_QCD/v liability explicit and mints op:qcd-scale.
%   - Flags the mechanism-level debts honestly: Goldstone/longitudinal
%     modes, rho = 1, WW unitarity (inherited phenomenologically, owed
%     structurally; adjacent to op:state-gate). Not minted as a new op --
%     folded into prose per the discussion.
%   - Notes the Koide-on-pole-masses puzzle as inherited, not resolved.
%   - CKM deliberately EXCLUDED from the accounting table (the delta-power
%     ansatz was not wired into the paper; claiming it here would
%     overstate).
%   - Literature positioning (Froggatt--Nielsen, warped-localisation kin)
%     deliberately omitted from this subsection; belongs in related-work
%     if wanted.
%
%  CROSS-REFS used (all verified to exist): sec:masses, sec:generations,
%  sec:confinement, sec:predictions, op:sector-scales, op:vev,
%  op:higgs-lambda, op:associator-geodesic, op:confinement-potential,
%  op:state-gate, sec:fanout-boundary (in hopf_justification -- verify
%  present in your tree). Citation keys existing: szangolies2025standardmodel.
%  NEW citation keys (RIS in higgs_coupling_refs_2026-07-03.ris):
%  atlas2022higgs, cms2022higgs, cepeda2019hllhc.
%  British English throughout.
% =====================================================================

\subsection{The associator debt and the Higgs mechanism}
\label{sec:associator-higgs}

Two accounts of one quantity have now been given, and they must be shown to be one account. The
Standard Model's account is that every elementary mass arises from coupling to the Higgs condensate:
a fermion mass is $m_f = y_f v/\sqrt2$, gauge-invariantly generated, with the Yukawa coupling $y_f$
a free parameter --- nine of them across the charged fermions, unexplained, constituting the flavour
puzzle. This framework's account is that mass is associator debt (\S\ref{sec:masses}): existence from
the associator, ratios from the triality angle and the dilution rule, and one free dimensional scale.
A reader is entitled to ask which of these the electron's mass actually is, and the answer this
subsection defends is: both, because they are the same statement. The Standard Model explains how
mass generation can be gauge-invariant; it does not explain what the masses are. The framework
derives what the masses are, and in doing so recovers --- rather than rivals --- the mechanism.

The identification is forced by the form of the results already in hand. Every mass formula in this
section and in \S\ref{sec:masses} has the shape $m = (\text{dimensionless geometric factor}) \times
v$: the scale was identified there as the Higgs vacuum value, and the $\sqrt2$ doubling norm as its
algebraic signature. The Standard Model writes the same relation as $m_f = y_f\,v/\sqrt2$ and calls
the dimensionless factor a Yukawa coupling. The framework's geometric factors therefore \emph{are}
the Yukawa couplings, $\sqrt2$-normalisation included, and the preceding subsection has already
exhibited one: $y_t = 1$ is a derived Yukawa coupling, stated in the Standard Model's own notation.
In the same language, the Koide relation and the dilution rule of \S\ref{sec:masses} are a
zero-parameter derivation of the Yukawa \emph{ratios} within each charge sector, and the sector-scale
problem (Open Problem~\ref{op:sector-scales}) acquires its proper name: it is exactly the
Yukawa-normalisation problem --- derive the lepton and down-sector normalisations from $v$ as
$y_t = 1$ already does for the up sector. The identification mints no new debt; it names an existing
one correctly.

The mechanism itself is recovered, not merely matched numerically. The associator debt is measured
\emph{relative to the selected complex direction}: it is the obstruction met in connecting the two
handedness resolutions of a frustrated pair (\S\ref{sec:generations}) through the preferred
$\mathbb{C} \subset \mathbb{O}$, which is the framework's rendering of the fact that a bare fermion
mass term is forbidden and the connection must be mediated --- the same shape as the Yukawa vertex
joining the doublet to the singlet through the Higgs field. The Higgs boson has already been
identified above as the dynamical promotion of that preferred direction
\autocite{szangolies2025standardmodel}, with vacuum value $v$ the magnitude of the selection. It
follows that in the symmetric phase there is no selected direction against which any debt can be
measured, and every fermion is massless: this is precisely electroweak symmetry breaking, recovered
as a statement about the selection rather than assumed. The question the framework defers --- why a
measurement selects exactly one imaginary direction (\S\ref{sec:fanout-boundary}) --- is thereby
revealed to be the question the Standard Model also cannot answer, in its own words: why does the
electroweak vacuum break the symmetry at all. One open question, two vocabularies.

The identification has an immediate consequence that experiment has already tested. The geometric
factors contain no reference to the scale --- they are built from fibre dimensions, the triality
angle and the Cayley--Dickson depth --- so the debt is exactly linear in $v$, and the coupling of the
physical Higgs to a fermion is $\partial m_f/\partial v = m_f/v$: coupling proportional to mass, the
Standard Model's signature relation, obtained here as a one-line theorem of linearity. This is
precisely the pattern the LHC has established --- the third-generation couplings observed by both
experiments, evidence at the muon, and agreement with mass-proportionality across three orders of
magnitude in mass \autocite{atlas2022higgs,cms2022higgs}. On any reading in which mass were intrinsic
geometry unrelated to the condensate, that measured proportionality would be an unexplained
coincidence; on the identification it could not have been otherwise. The gauge-boson couplings
$2m_V^2/v$ follow identically, the $W$ and $Z$ masses being linear in $v$ through the
vacuum-engagement of the classification above; and the light neutrinos are the one correct exception,
their masses quadratic in $v$ through the seesaw because the a-chiral partner's Majorana scale is not
Higgs-generated at all --- a gauge singlet needs no condensate, exactly as in the standard seesaw
(\S\ref{sec:predictions}).

Because the framework derives the Yukawa couplings rather than modifying the mechanism, its
effective theory beneath the matching scale is the Standard Model with these coupling values and
nothing else, and the prediction is correspondingly rigid: every Higgs coupling-modifier is unity,
$\kappa_f = \kappa_V = 1$, at every attainable precision, forever. The high-luminosity programme is
expected to sharpen the third-generation determinations to the few-per-cent level and to reach the
second generation \autocite{cepeda2019hllhc}; each must land at the Standard-Model point, and a
single established deviation anywhere refutes the framework outright. Where most extensions of the
Standard Model predict departures somewhere, this framework forbids them everywhere; the desert below
the matching scale and the exactness of the couplings are two faces of the same claim, and the
prediction is collected with the other permanent nulls in \S\ref{sec:predictions}.

Honesty requires an equally plain statement of what the identification does not supply. The Higgs
mechanism delivers, beyond the masses, the Goldstone bosons as the longitudinal modes of the $W$ and
$Z$, the tree-level relation $\rho = m_W^2/(m_Z^2\cos^2\theta_W) = 1$, and the unitarisation of
$WW$ scattering by Higgs exchange; through the identification the framework inherits these
phenomenologically, since its effective theory is the Standard Model, but exhibiting them within the
framework's own structures is open and sits with the state-to-gate correspondence
(Open Problem~\ref{op:state-gate}). The unification of \S\ref{sec:confinement} --- all mass is
associator-geodesic length --- likewise obliges the $W$ and $Z$ masses, derived above by
vacuum-engagement, to be exhibited as geodesic lengths like every other mass; the identification
suggests how, engagement with the selection being the debt, but this is a structural claim, not a
proof. And the radiative structure is only partly understood: the geometric values are tree-level,
and the residuals carry the size and sign of loop dressing --- $y_t$ high by $0.9\%$ against the pole
mass, $m_H$ low by $1.7\%$, the latter already posed as a running consistency check in
Open Problem~\ref{op:higgs-lambda} --- but the survival of the Koide relation on \emph{pole} masses
to one part in $10^5$, radiative dressing included, is a puzzle this framework inherits from Koide's
own literature and does not resolve: either the geometry constrains the dressed quantities, a strong
claim we do not make, or the robustness is unexplained. We flag it rather than paper over it.

One further bill comes due here, and it is a direct consequence of commitments made elsewhere in
this paper. Under the unification of \S\ref{sec:confinement}, the hadronic scale is not exempt from
the accounting: if all mass is associator-geodesic length and the framework carries exactly one free
dimensional parameter, then the ratio of the hadronic scale to the electroweak scale --- of order
$10^{-3}$, the number that makes the proton weigh a GeV rather than a fraction of $v$ --- is a
dimensionless quantity the framework is committed to deriving. The Standard Model treats the two
scales as independent inputs; this framework, by its own accounting, may not.

\begin{openproblem}[The hadronic scale]\label{op:qcd-scale}
Derive the ratio of the hadronic scale to the electroweak scale, $\Lambda_{\mathrm{QCD}}/v$ of order
$10^{-3}$ --- equivalently, the proton mass in units of $v$ --- as the associator-geodesic length of
the closed colour network (Open Problems~\ref{op:associator-geodesic}
and~\ref{op:confinement-potential}) in units of the vacuum value. Under the single-parameter claim of
\S\ref{sec:masses} this ratio is not an independent input but an owed prediction, and it is stated
here so that the claim's full price is visible.
\end{openproblem}

The accounting can now be drawn up. Of the eleven numbers of the Higgs sector proper --- the nine
Yukawa couplings, the quartic, and the vacuum value --- the framework derives eight, leaves two open,
and holds one structurally free. (The CKM sector is addressed only partially in this paper and is
deliberately excluded from this accounting.)

\begin{center}
\begin{tabular}{lll}
\hline
Higgs-sector quantity & Standard Model & this framework \\
\hline
quartic $\lambda$ & free & $\tfrac18$ ($-3.4\%$; Open Problem~\ref{op:higgs-lambda}) \\
vacuum value $v$ & measured input & the single free scale (Open Problem~\ref{op:vev}) \\
top Yukawa $y_t$ & free & $1$ ($+0.9\%$ against the pole mass) \\
within-sector Yukawa ratios (six) & free & $\delta = \tfrac29$ and dilution: leptons to $10^{-5}$, \\
 & & quarks $\lesssim 1\%$ ($u$ at $+29\%$, flagged) \\
cross-sector normalisations (two) & free & open (Open Problem~\ref{op:sector-scales}) \\
coupling modifiers $\kappa_f$, $\kappa_V$ & $\equiv 1$ by construction & $= 1$ by linearity: exact prediction \\
\hline
\end{tabular}
\end{center}

The subsection's claim, in one sentence: the Higgs mechanism is how nature makes mass generation
gauge-invariant, the associator debt is what the masses are, and the two meet in the Yukawa
couplings --- which the Standard Model measures, and this framework computes.

% =====================================================================
%  COMPANION SNIPPET -- paste into sec:predictions, in the permanent-null
%  ("null-forever") family, alongside FCNC / no proton decay / desert.
%  One paragraph; uses the new label sec:associator-higgs and the key
%  cepeda2019hllhc.
% =====================================================================
%
% \paragraph{Exact Standard-Model Higgs couplings.}
% Because the framework derives the Yukawa couplings rather than modifying
% the mechanism (\S\ref{sec:associator-higgs}), every Higgs
% coupling-modifier is exactly unity: $\kappa_f = \kappa_V = 1$, at every
% attainable precision, permanently. The high-luminosity programme's
% few-per-cent third-generation determinations and the coming
% second-generation measurements at the muon and charm
% \autocite{cepeda2019hllhc} must all land at the Standard-Model point; a
% single established deviation anywhere refutes the framework outright.
% Where most extensions of the Standard Model predict departures somewhere,
% this framework forbids them everywhere --- the desert below the matching
% scale and the exactness of the couplings are two faces of one claim.
% =====================================================================
